\documentclass[11pt]{article}
\usepackage{fontspec,amsmath,amssymb,amsthm,hyperref}
\usepackage[margin=1in]{geometry}
\setlength{\emergencystretch}{3em}
\title{Project splitting transport to the original parameter center}
\author{LangevinLean mathematical exposition}
\date{}
\begin{document}\maketitle
\noindent Source module: \texttt{ProximalLangevin/SplittingProjection.lean}. Compile this proof document with XeLaTeX or LuaLaTeX. The accompanying HTML places the complete Lean declarations beside every proof.


Let \(\rho_n\) be the actual PIPGLA law with any positive proximal parameter \(\lambda\), and compare it with the Gibbs target for \(V_\eta=g_1+g_2^\eta\), \(\eta>0\). This module does not assume the algorithm already has a transport convergence rate.


\section{Establish state and projected second moments}


The law equals the pushforward of the constructed \(L^2\) chain. The map-measure criterion then proves \(z\mapsto z\in L^2(\rho_n)\). Apply the bounded linear parameter map and subtract a constant to get \(\theta(z)-a\in L^2(\rho_n)\) for every center \(a\).



\noindent Lean declarations: \texttt{pipglaLaws\_state\_memLp}, \texttt{pipglaLaws\_parameter\_memLp}.

\section{Combine transport, equilibrium concentration and center displacement}


Sharp marginal concentration gives target RMS at most \(\sqrt{d_\theta/(Nm)}\), and the actual uniform potential gap gives \(\|\theta_U-\theta_\eta\|\le\sqrt{\eta K^2/m}\). Nonexpansive parameter projection and the center triangle inequality yield

\[\operatorname{RMS}_{\rho_n}(\theta_U)\le\sqrt{\eta K^2/m}+\sqrt{d_\theta/(Nm)}+W_2(\rho_n,\Pi_{N,\eta}).\]

State and target moment bounds establish finite transport before the real-valued projection inequality is used.



\noindent Lean declarations: \texttt{pipgla\_parameter\_rms\_original\_le\_transport}.

\end{document}
